\documentclass[10pt,a4paper]{article}
\usepackage[margin=2.1cm]{geometry}
\usepackage{amsmath,amssymb,bm}
\usepackage{graphicx,booktabs,caption}
\usepackage[hidelinks]{hyperref}
\usepackage{lmodern}
\usepackage[T1]{fontenc}
\usepackage{microtype}
\usepackage{xcolor}
\usepackage{enumitem}
\captionsetup{font=small,labelfont=bf}
\newcommand{\NL}{63}
\newcommand{\Mtrain}{29\,312}
\newcommand{\Mtest}{7\,401}
\newcommand{\Mnovel}{7401}
\newcommand{\Mcopies}{0}
\newcommand{\PctCopies}{0}
\newcommand{\ThetaTrain}{0.312}
\newcommand{\MeffTrain}{3590}
\newcommand{\ThetaTest}{0.303}
\newcommand{\MeffTest}{1280}
\newcommand{\MeffNovel}{1280}
\newcommand{\NfOne}{1\,323}
\newcommand{\NCij}{861\,273}
\newcommand{\Hbest}{128}
\newcommand{\dbest}{10}
\newcommand{\lambest}{10^{-3}}
\newcommand{\Nparams}{217\,728}
\newcommand{\BestEpoch}{668}
\newcommand{\ValNLPL}{82.18}
\newcommand{\NLPLtrain}{71.10}
\newcommand{\NLPLtest}{78.72}
\newcommand{\NLPLnovel}{78.72}
\newcommand{\NLPLgap}{7.6}
\newcommand{\NLPLgapPct}{11}
\newcommand{\NLPLtestSite}{1.25}
\newcommand{\PPVL}{0.83}
\newcommand{\PPVLhalf}{1.00}
\newcommand{\PPVbase}{0.21}
\newcommand{\PCone}{11.3}
\newcommand{\PCtwo}{4.3}
\newcommand{\Etest}{-131.7}
\newcommand{\Etrain}{-134.3}
\newcommand{\Erandom}{-0.0}
\newcommand{\HamTestTest}{0.60}
\newcommand{\HamTestNN}{0.12}
\newcommand{\HamNovelTrain}{0.18}
\newcommand{\HamTestTrain}{0.18}
\newcommand{\rCtrain}{0.944}
\newcommand{\rCprof}{-0.005}
\newcommand{\sCtrain}{0.88}
\newcommand{\Crel}{0.94}
\newcommand{\sCdisTrain}{0.94}
\newcommand{\sCdisOne}{0.83}
\newcommand{\sCdisStar}{0.91}
\newcommand{\rCnovelFull}{1.000}
\newcommand{\BetaStar}{1.05}
\newcommand{\TauOne}{6.5}
\newcommand{\TauPOne}{3.8}
\newcommand{\BurnOne}{500}
\newcommand{\NgenOne}{10\,000}
\newcommand{\NcolOne}{0}
\newcommand{\PctColOne}{0.0}
\newcommand{\StudyColNatOne}{0.0}
\newcommand{\StudyColRndOne}{0.0}
\newcommand{\HamEndOne}{0.66}
\newcommand{\SweepsNinetyOne}{11}
\newcommand{\EendNatOne}{-126.6}
\newcommand{\EendRndOne}{-126.6}
\newcommand{\EgenOne}{-125.9}
\newcommand{\rFOne}{0.990}
\newcommand{\rCOne}{0.871}
\newcommand{\sCOne}{0.78}
\newcommand{\rCnovelOne}{0.871}
\newcommand{\HamGGOne}{0.68}
\newcommand{\HamGTOne}{0.44}
\newcommand{\HamGTrOne}{0.42}
\newcommand{\HamGTrMinOne}{0.32}
\newcommand{\UniqOne}{100.0}
\newcommand{\StudySweeps}{2000}
\newcommand{\WPCaOne}{0.26}
\newcommand{\WPCbOne}{0.34}
\newcommand{\WEOne}{5.9}
\newcommand{\TauStar}{5.6}
\newcommand{\TauPStar}{3.4}
\newcommand{\BurnStar}{500}
\newcommand{\NgenStar}{10\,000}
\newcommand{\NcolStar}{0}
\newcommand{\PctColStar}{0.0}
\newcommand{\StudyColNatStar}{0.0}
\newcommand{\StudyColRndStar}{0.0}
\newcommand{\HamEndStar}{0.64}
\newcommand{\SweepsNinetyStar}{17}
\newcommand{\EendNatStar}{-134.6}
\newcommand{\EendRndStar}{-134.6}
\newcommand{\EgenStar}{-133.8}
\newcommand{\rFStar}{0.992}
\newcommand{\rCStar}{0.859}
\newcommand{\sCStar}{0.85}
\newcommand{\rCnovelStar}{0.859}
\newcommand{\HamGGStar}{0.65}
\newcommand{\HamGTStar}{0.41}
\newcommand{\HamGTrStar}{0.39}
\newcommand{\HamGTrMinStar}{0.29}
\newcommand{\UniqStar}{100.0}
\newcommand{\WPCaStar}{0.30}
\newcommand{\WPCbStar}{0.31}
\newcommand{\WEStar}{2.1}
\newcommand{\WPCaTrain}{1.23}
\newcommand{\WPCbTrain}{0.29}
\newcommand{\WPCaProf}{1.05}
\newcommand{\WPCbProf}{0.66}
\newcommand{\WEtrain}{2.7}
\newcommand{\WEprof}{39.5}
\newcommand{\Eunw}{-146.2}
\newcommand{\ClTest}{14.1}
\newcommand{\ClTrain}{12.9}
\newcommand{\ClOne}{5.1}
\newcommand{\ClStar}{3.2}
\newcommand{\Kclust}{3}
\newcommand{\NcovClust}{2}
\newcommand{\CovWeight}{85.9}
\newcommand{\NtestCov}{6\,761}
\newcommand{\MeffTestCov}{1102}
\newcommand{\rCtrainCov}{0.941}
\newcommand{\rFcovOne}{0.991}
\newcommand{\rCcovOne}{0.918}
\newcommand{\sCcovOne}{0.82}
\newcommand{\rCcovTopOne}{0.961}
\newcommand{\rFcovStar}{0.996}
\newcommand{\rCcovStar}{0.918}
\newcommand{\sCcovStar}{0.91}
\newcommand{\rCcovTopStar}{0.961}
\newcommand{\NsuperSet}{36\,713}
\newcommand{\OrigCopies}{78.8}
\newcommand{\OrigNN}{0.044}
\newcommand{\SplitThresh}{0.2}
\newcommand{\SplitIdent}{80}
\newcommand{\NclustTot}{10\,942}
\newcommand{\NsingleClust}{6\,767}
\newcommand{\NtrainClean}{29\,312}
\newcommand{\NtestClean}{7\,401}
\newcommand{\NdropClean}{0}
\newcommand{\CleanNN}{0.179}
\newcommand{\CleanNNmin}{0.016}
\newcommand{\CleanNNminSub}{1}
\newcommand{\CleanNNmed}{0.175}
\newcommand{\CleanNNqfive}{0.032}
\newcommand{\CleanPctNear}{61}
\newcommand{\CleanDropRule}{exact}
\newcommand{\CmpMtrain}{14\,502}
\newcommand{\CmpMtest}{3\,739}
\newcommand{\CmpPctCopies}{79}
\newcommand{\CmpValNLPL}{85.61}
\newcommand{\CmpNLPLtest}{75.70}
\newcommand{\CmpNLPLnovel}{77.13}
\newcommand{\CmprFOne}{0.990}
\newcommand{\CmprCOne}{0.865}
\newcommand{\CmprCtrain}{0.960}
\newcommand{\CmpHamGTrOne}{0.42}
\newcommand{\CmpPPVLhalf}{0.97}
\newcommand{\CmprCcovOne}{0.921}
\newcommand{\MissCl}{C0}
\newcommand{\EmissTest}{-119.5}
\newcommand{\EmissTrain}{-122.6}
\newcommand{\EcovTest}{-134.0/-133.5}
\newcommand{\EcovTrain}{-134.5/-137.2}
\newcommand{\EgapMiss}{14}
\newcommand{\PLLmissTrain}{-71.8}
\newcommand{\PLLcovTrain}{-71.1/-71.0}
\newcommand{\dFmiss}{+1.01}
\newcommand{\dFmissFac}{0.36}


\newcommand{\cost}[1]{\textsf{\footnotesize #1}}
\newcommand{\PC}{\mathrm{PC}}
\newcommand{\PCA}{\mathrm{PCA}}

\title{\vspace{-1.2cm}Conditioning generation on PCA coordinates:\\ post-mortem and candidate methods}
\author{AttentionDCA / J-domains --- proposal for the next experiment}
\date{\today}

\begin{document}
\maketitle

\begin{abstract}
\noindent
An earlier study appended the discretised PCA coordinates to each sequence as two extra sites and
learned a joint coupling tensor over sequence and label, steering generation with a separate
inverse temperature $\beta_{\PCA}$. It works only partially: steering requires
$\beta_{\PCA}\gtrsim10$, and beyond that the samples stop being J-domains. This note diagnoses
\emph{why} that happens, then sets out four ways to do conditional generation properly. The
central observation is that a PCA coordinate is an \emph{exactly linear} function of the one-hot
sequence, so conditioning on it needs no learning at all: the constraint can be written in closed
form as a Potts term and added to the existing model. Nothing here is implemented yet; a common
evaluation protocol is given at the end. A secondary section covers the separate question of the
unconditional model's mode weights.
\end{abstract}

%======================================================================
\section{Two goals that are often confused}
%======================================================================
\begin{description}[leftmargin=1.2em,itemsep=2pt]
\item[\textbf{(G1) Conditional generation.}] ``Give me sequences at this place in PCA space.''
  Success means the samples land near the target \emph{and} remain valid J-domains. This is the
  goal of the earlier study.
\item[\textbf{(G2) Correct mode weights.}] ``When I sample freely, each sub-family should appear
  with its natural frequency.'' Our current model fails this for one cluster: \MissCl\ carries
  \ClTest\% of the natural weight but receives \ClOne\% of the samples ($\Delta F=\dFmiss$ nats).
\end{description}
They are different problems. G1 is solved by adding a constraint at sampling time; G2 is a
property of the learned energy function and is not fixed by conditioning --- though, usefully, a
good solution to G1 gives a \emph{measurement} tool for G2 (\S\ref{sec:fe}). Sections
\ref{sec:post}--\ref{sec:cond} address G1; \S\ref{sec:uncond} summarises G2.

%======================================================================
\section{Post-mortem: why appending discretised coordinates underperforms}\label{sec:post}
%======================================================================
In the earlier work the $35\times35$ discretised PCA grid was appended as two extra sites with an
enlarged alphabet ($\bar q=35$), giving a joint tensor $J$ of shape
$(\bar q,\bar q,L{+}2,L{+}2)$, and sampling used
$P(a_i\mid\bm a_{-i},\bm z)\propto\exp\big(\sum_{j\le L}J_{ij}(a_i,a_j)+\beta_{\PCA}\sum_{j>L}
J_{ij}(a_i,z_j)\big)$.
Four structural problems explain the reported behaviour, and none of them is a tuning issue.

\begin{enumerate}[leftmargin=*,itemsep=3pt]
\item \textbf{The label is outvoted $2$ to $L$.} In the conditional for site $i$ the label
      contributes $2$ of $L+2=65$ coupling terms, about $3\%$ of the local field. That is exactly
      why steering only begins around $\beta_{\PCA}\approx10$. But $\beta_{\PCA}$ multiplies
      \emph{learned} couplings, whose directions in sequence space are arbitrary wherever the data
      were thin; amplifying them by $10$--$50$ pushes sequences along those arbitrary directions,
      which is the observed ``shifts too far right, loses the HPD motif''. There is no
      $\beta_{\PCA}$ that both steers and preserves the protein, because the two effects are
      driven by the same amplified parameters.
\item \textbf{One-hot bins throw away the geometry of PCA space.} Bins $20$ and $21$ are adjacent
      in $\PC_1$ but are unrelated symbols in a Potts alphabet: nothing ties their columns
      together. With $35\times35=1225$ cells and ${\sim}10^4$ training sequences, most cells hold a
      handful of sequences or none, so the label couplings are estimated from almost no data and
      cannot interpolate. This is the direct explanation of the coordinate-generation figures:
      conditioning on a true training sequence and sampling its coordinates puts mass on one or
      two spurious cells (Model 1/2) or is nearly uniform (Model 3), instead of recovering the
      true cell.
\item \textbf{The label is a deterministic function of the sequence.} $\bm z=f(\bm x)$ exactly, so
      the true $p(\bm z\mid\bm x)$ is a delta. A handful of pairwise terms to a 35-state variable
      cannot represent a delta, so the fit is forced to be wrong; the model is being asked to
      learn something it cannot express, and the coordinate-generation test measures precisely
      that failure.
\item \textbf{It is unnecessary.} The quantity being conditioned on is already an exact linear
      function of the model's own variables (\S\ref{sec:cond}), so learning it from data only
      introduces estimation error where a closed form exists.
\end{enumerate}

%======================================================================
\section{The key fact: $\PC$ is linear in the one-hot sequence}\label{sec:cond}
%======================================================================
Let $\bm x\in\{0,1\}^{Lq}$ be the one-hot encoding, $x_{ia}=1$ iff site $i$ carries letter $a$.
The PCA used throughout this project is fitted on exactly these vectors, so for component $c$
\begin{equation}
\PC_c(\bm x)=\sum_{i,a}V_{c,ia}\,x_{ia}-m_c,\qquad m_c=\sum_{i,a}V_{c,ia}\,\mu_{ia},
\label{eq:lin}
\end{equation}
with $V$ the loadings and $\mu$ the training mean. This is not an approximation: evaluating
\eqref{eq:lin} against \texttt{sklearn}'s \texttt{transform} on our test set reproduces it to
$0.0$ (machine zero).

Therefore a harmonic restraint towards a target $\bm t$ in PCA space,
\begin{equation}
E_\lambda(\bm x)=E(\bm x)+\frac{\lambda}{2}\sum_{c}\big(\PC_c(\bm x)-t_c\big)^2 ,
\label{eq:restraint}
\end{equation}
is \emph{exactly a Potts energy}: expanding the square gives a pairwise part and a field part,
\begin{equation}
\Delta J_{ia,jb}=-\lambda\sum_c V_{c,ia}V_{c,jb},\qquad
\Delta h_{ia}=\lambda\sum_c (t_c+m_c)\,V_{c,ia},
\end{equation}
i.e.\ a rank-$C$ correction to the coupling matrix plus a field. No training, no new parameters,
no $\beta_{\PCA}$ to tune blindly --- $\lambda$ is a restraint stiffness with a clear meaning, as
in umbrella sampling.

\paragraph{Why this behaves better than the learned version.} Increasing $\lambda$ constrains the
sequence only along the $C$ directions $V_c$; every direction orthogonal to them stays governed by
the protein energy $E(\bm x)$. The failure mode of the earlier approach --- amplifying arbitrary
learned directions until the protein breaks --- cannot occur, because the constrained direction is
the PCA direction itself, by construction. Quality still degrades as $\lambda\to\infty$ (the
constraint eventually dominates), but gracefully and for an understandable reason.

\paragraph{Scale, for intuition.} On our clean split, $\PC_1$ spans $6.9$ across the family with
standard deviation $2.06$, while a single substitution moves it by $0.125$ on average (at most
$0.42$). Reaching a distant target therefore requires tens of coordinated substitutions, which is
why a weak restraint does nothing and why the restraint must be applied \emph{during} sampling
rather than as a post-hoc filter.

\subsection*{C1 --- Analytic harmonic restraint (recommended)}
Implement \eqref{eq:restraint} directly in the existing sampler. Each chain keeps a running
$\PC_c$; changing site $i$ from $a$ to $a'$ shifts it by $V_{c,ia'}-V_{c,ia}$, so the conditional
logits cost $O(q\,C)$ extra per site update --- negligible.

\noindent\cost{Cost:} ${\sim}30$ lines in \texttt{sampler.py}, no retraining, minutes to run.
\cost{Expected:} samples concentrated near $\bm t$ with a controlled trade-off; a monotone,
interpretable $\lambda$ curve replacing the non-monotone $\beta_{\PCA}$ behaviour.
\cost{Risks:} a harmonic restraint targets a \emph{point}; for a region, restrain to the region's
centre with a soft $\lambda$, or use a flat-bottomed restraint (zero penalty inside the target
cell, quadratic outside). Also, $\PC$ is a 2-D summary explaining only
$\PCone\%+\PCtwo\%$ of the variance, so ``hit the target'' never means ``is the right protein''
--- sequence-quality metrics must be reported alongside.

\subsection*{C2 --- Exact conditional by rejection (ground truth)}
Sample the unconditioned model, keep only the sequences whose coordinates fall in the target cell.
This yields \emph{exactly} $p(\bm x\mid \bm z\in\text{cell})$ with no approximation.

\noindent\cost{Cost:} 20 lines, no retraining. \cost{Expected:} the reference answer every other
method should be compared against --- for the first time we would know what the conditional
distribution actually is. \cost{Risks:} efficiency falls with the natural density of the target;
for a cell in the dense part of the family it is cheap, for the far tail it is hopeless. This is
the method's honest limitation, and it is exactly why C1/C3 exist --- but it makes C2 the right
validation for \emph{reachable} targets, including \MissCl.

\subsection*{C3 --- Learned conditional field (if a learned model is wanted)}
Keep the coordinates \emph{continuous} and let them bias every site through a field rather than
through two coupling columns:
\begin{equation}
p(x_i=a\mid\bm x_{-i},\bm z)\propto\exp\Big(\sum_{j\neq i}J_{ij}(a,x_j)+\bm w_{ia}\!\cdot\!\phi(\bm z)\Big),
\qquad \phi(\bm z)=(1,z_1,z_2,z_1^2,z_2^2,z_1z_2),
\end{equation}
trained by conditional pseudo-likelihood on $(\bm x,\bm z=f(\bm x))$ pairs. Parameters:
$|\phi|\,Lq\approx8\,000$, a $4\%$ increase. This repairs problems 1 and 2 of \S\ref{sec:post}:
the label now reaches \emph{every} site, and smoothness in $\bm z$ is built into the
parametrisation so neighbouring coordinates share statistical strength.

\noindent\cost{Cost:} moderate --- conditional field in model, trainer and sampler; 1 training run
($\approx40$ min). \cost{Expected:} steering without a large temperature, and a model that can also
answer $p(\bm z\mid\bm x)$ sensibly. \cost{Risks:} still learns what C1 knows exactly, so it is
only worth it if the aim is a \emph{learned} conditional model (e.g.\ to later condition on
something that is \emph{not} a linear function of the sequence --- phenotype, host, temperature).
That is the real argument for C3, and it is a good one.

\subsection*{C4 --- Repair the label-as-sites model (if continuity with the old work matters)}
Two targeted fixes: (i) give the label an ordinal/smooth encoding --- e.g.\ replace the $35$-state
one-hot bin by a small set of basis features, or tie neighbouring bins with a smoothness penalty
$\sum_z\|J(\cdot,z{+}1)-J(\cdot,z)\|^2$; (ii) rebalance the $2$-vs-$L$ asymmetry by replicating the
label over $R$ sites or by scaling the label block at \emph{training} time rather than inflating
$\beta_{\PCA}$ at sampling time.

\noindent\cost{Cost:} 1 training run plus the smoothness term. \cost{Expected:} coordinate
generation that recovers the true cell, which is the clean diagnostic the old report found broken.
\cost{Risks:} keeps the discretisation that caused the trouble; strictly dominated by C1 for this
particular conditioning variable.

\begin{table}[h]\centering\small
\begin{tabular}{llccl}
\toprule
 & learning needed & effort & $\beta_{\PCA}$ pathology & what it gives \\
\midrule
C1 harmonic restraint & none (closed form) & low    & no  & controlled steering, free-energy maps \\
C2 rejection          & none               & trivial& no  & the exact conditional, as ground truth \\
C3 conditional field  & yes                & medium & no  & a learned conditional, extensible to other labels \\
C4 repaired labels    & yes                & medium & reduced & continuity with the previous study \\
\bottomrule
\end{tabular}
\caption{Conditional-generation options. C1 and C2 require no retraining.}
\end{table}

%======================================================================
\section{A bonus: the restraint measures the model's error everywhere}\label{sec:fe}
%======================================================================
Running C1 at a grid of targets with a moderate $\lambda$ is umbrella sampling along
$(\PC_1,\PC_2)$. Reweighting the windows (WHAM, or simply the known biasing potential) gives the
model's free-energy surface $F_{\text{model}}(\bm z)$, which can be compared directly with the
natural one, $-\ln\rho_{\text{nat}}(\bm z)$ estimated from the reweighted training sequences. The
difference is a \emph{map} of where the model misplaces probability --- a much richer object than
the three per-cluster numbers we currently report, and it would show whether the \MissCl\ deficit
is a uniform offset of that basin or a distortion of its shape. \cost{Cost: one C1
implementation plus a grid of runs.}

%======================================================================
\section{Separately: the unconditional mode weights (G2)}\label{sec:uncond}
%======================================================================
For completeness, the options for the unconditional problem, unchanged from the previous note:
\textbf{M1} per-cluster reweighting of the training set by $e^{\gamma\Delta F_k}$ (a control loop;
cheap, but \MissCl\ has only $M_\text{eff}=\MeffMiss$ against \Nparams\ parameters, so it risks
over-fitting); \textbf{M2} add explicit fields $h_i$ (cheap, a correctness fix in its own right);
\textbf{M3} an explicit mixture with a hand-set prior $\pi_k$ (correct by construction, but imposes
rather than learns the prior); \textbf{M4} likelihood-based fine-tuning by persistent contrastive
divergence, whose gradient $\langle\cdot\rangle_{\text{data}}-\langle\cdot\rangle_{\text{model}}$
contains exactly the term pseudo-likelihood omits (the principled fix, highest effort);
\textbf{M5} resampling generated pools to the natural cluster shares (a control, not a fix).

Note that C1/C2 give a cheap partial answer to G2 as well: generating a pool whose cluster
proportions are forced to the natural ones isolates how much of the full-set statistics gap is due
to mode weight alone. Since restricting to the covered clusters already lifts $r(C_{ij})$ from
\rCOne\ to \rCcovOne, we can predict the outcome and check it.

%======================================================================
\section{Evaluation protocol (fixed in advance)}
%======================================================================
For conditional generation, every method is reported on the same two axes, because the previous
study's difficulty was precisely that improving one destroys the other:
\begin{itemize}[leftmargin=*,itemsep=1pt]
\item \textbf{On-target accuracy:} mean and spread of $\|\PC(\bm x)-\bm t\|$ over the generated
      pool; fraction landing in the target cell.
\item \textbf{Sequence quality, at matched accuracy:} $r(f_i)$ and $r(C_{ij})$ against the
      \emph{natural sequences in the same region} (not the whole family --- the earlier report
      noted this distinction and it is the right one); retention of the HPD motif; mean distance
      to the nearest natural sequence; held-out energy of the samples.
\item \textbf{The trade-off curve itself:} sweep $\lambda$ (or $\beta_{\PCA}$) and plot quality
      against accuracy. A method is better if its curve dominates, which is a fairer comparison
      than any single operating point.
\item \textbf{Against ground truth:} wherever C2 is affordable, report the gap between each
      method's conditional and the exact one.
\item \textbf{Targets:} at least one dense target and one in \MissCl, since the sparse region is
      where methods differ.
\end{itemize}

\vspace{0.2cm}
\noindent\textbf{Recommendation.} \textbf{C1 + C2 together}: C2 fixes what the right answer is,
C1 provides a cheap and well-behaved way to reach it, and neither needs retraining, so both can be
done and evaluated quickly. Add \textbf{C3} if the goal is a learned conditional model that will
later be conditioned on labels which are \emph{not} functions of the sequence --- that is the case
where the analytic route is unavailable and the extra machinery earns its keep.

\end{document}
